% Build with: tectonic paper.tex
\documentclass[11pt,letterpaper]{article}
\usepackage[T1]{fontenc}
\usepackage{lmodern}
\usepackage[margin=1in,headheight=15pt]{geometry}
\usepackage{amsmath,amssymb,amsthm}
\usepackage{microtype,booktabs,enumitem,xcolor,fancyhdr,hyperref}
\usepackage[most]{tcolorbox}
\usepackage{needspace}
\definecolor{ink}{HTML}{203B55}
\definecolor{accent}{HTML}{176C65}
\definecolor{boxbg}{HTML}{F2F6F8}
\hypersetup{colorlinks=true,linkcolor=ink,urlcolor=accent,citecolor=ink,
 pdftitle={Depth-Two Threshold Circuits for Andreev's Function: Sharper Bounds},
 pdfauthor={Christian Reitwiessner}}
\pagestyle{fancy}
\fancyhf{}
\fancyhead[L]{\small\color{ink}Threshold circuits for Andreev's function}
\fancyhead[R]{\small\thepage}
\renewcommand{\headrulewidth}{0.3pt}
\setlength{\parindent}{0pt}
\setlength{\parskip}{0.55em}
\setlength{\emergencystretch}{2em}
\setlist[itemize]{leftmargin=1.5em,itemsep=0.2em,topsep=0.2em}
\setlist[enumerate]{leftmargin=1.8em,itemsep=0.2em,topsep=0.2em}
\newtheorem{theorem}{Theorem}
\newtheorem{lemma}{Lemma}
\newtheorem{corollary}{Corollary}
\newtheorem{proposition}{Proposition}
\newcommand{\bits}{\{0,1\}}
\newcommand{\ind}[1]{\bigl[\,#1\,\bigr]}
\newcommand{\sect}[1]{\section{\color{ink}#1}}
\newcommand{\subsect}[1]{\subsection*{\color{ink}#1}}
\newcommand{\smalltitle}[1]{\textbf{\color{ink}#1}\quad}
\newcommand{\THR}{\mathsf{THR}}
\newcommand{\MAJ}{\mathsf{MAJ}}
\newtcolorbox{keybox}[1][]{colback=boxbg,colframe=ink!40,boxrule=0.4pt,
  arc=2pt,left=8pt,right=8pt,top=5pt,bottom=5pt,#1}
\begin{document}
\thispagestyle{empty}
{\color{accent}\small RESEARCH NOTE\par}
\vspace{0.4em}
{\raggedright\color{ink}\LARGE\bfseries Depth-Two Threshold Circuits for\\[0.15em]
Andreev's Function: Sharper Bounds\par}
\vspace{0.5em}
{\large Christian Reitwiessner\par}
{\small November 2026\par}
\vspace{0.8em}

\begin{keybox}
\textbf{Summary.} Andreev's function on $n$ input bits reads a bit from a
table at an address formed by block parities. For depth-two threshold
circuits ($\THR\circ\THR$) we show:
\begin{itemize}
\item an exact circuit with $\tfrac34 n^2+1$ gates (previously
      $O(n^3/\log n)$, Kane and Williams~\cite{KW});
\item matching average-case bounds: agreeing with the function on a
      $\tfrac12+\varepsilon$ fraction of inputs takes
      $\tilde\Theta(\varepsilon^2 n^{3/2})$ gates (previously only the lower
      bound $\tilde\Omega(\varepsilon^3 n^{3/2})$ was known);
\item a worst-case lower bound of $\Omega(n^{3/2}/\log^2 n)$ gates and
      $\Omega(n^{5/2}/\log^{5/2}n)$ wires (previously $\log^3$ and
      $\log^{7/2}$).
\end{itemize}
For depth three we improve the top fan-in bound for Andreev's function, and
the gate and wire bounds for the harder function $B_n$ of Kane and Williams;
we also give the first average-case bound for $B_n$.
\end{keybox}

\sect{Setting and results}

\smalltitle{Andreev's function.}
Fix integers $k,m\ge1$ and put $N=2^k$. The input consists of
\begin{itemize}
\item a \emph{table} of $N$ bits $x_c$, one for each address $c\in\bits^k$, and
\item $k$ \emph{address blocks} of $m$ bits; block $i$ has \emph{weight}
      $s_i$, the number of ones in it.
\end{itemize}
The address $b\in\bits^k$ has coordinates $b_i=s_i\bmod2$, and
\[
A(x,a)=x_b .
\]
The input length is $n=N+km$. The \emph{balanced} case is $km=N$, so $n=2N$
and $k=\Theta(\log n)$. All asymptotic statements refer to the balanced case.

\smalltitle{Threshold circuits.}
For a statement $P$, the bracket $\ind{P}$ is $1$ if $P$ holds and $0$
otherwise. A \emph{threshold gate} with inputs $z_1,\dots,z_r$, real weights
$w_j$ and threshold $\tau$ computes $\ind{\sum_j w_jz_j\ge\tau}$. In a
\emph{depth-two circuit} the bottom gates read input bits and a single top
gate reads bottom gates and input bits. Its \emph{size} is its number of
gates; its \emph{wires} are its nonzero weights. A circuit
\emph{$\varepsilon$-agrees} with $A$ if it outputs $A(x,a)$ on at least a
$\tfrac12+\varepsilon$ fraction of all inputs.

The lower bounds hold for $\varepsilon\ge C\log n/\sqrt n$; below that,
a circuit with $O(1)$ gates already agrees with $A$ on a
$\tfrac12+\Omega(1/\sqrt n)$ fraction (majority of the table).
The remaining gap is the worst case: between
$n^{3/2}/\log^2n$ and $\tfrac34 n^2$.

\needspace{8\baselineskip}
\smalltitle{Results and the state of the art.}
Stronger lower bounds against $\THR\circ\THR$ are known for functions in
the larger class $\mathsf{E}^{\mathsf{NP}}$, proved by the algorithmic method.
For functions in $\mathsf P$, which includes Andreev's function and $B_n$, the
restriction method used here gives the strongest gate bounds we know of.
\begin{center}
\small
\renewcommand{\arraystretch}{1.25}
\begin{tabular}{@{}lllll@{}}
\toprule
Function & Bound & Best previously known & This paper & \\
\midrule
\multicolumn{5}{@{}l}{\emph{Functions in $\mathsf P$}}\\
Andreev, exact & gates, upper & $O(n^3/\log n)$ \cite{KW} & $\tfrac34n^2+1$ & Thm.~\ref{thm:exact}\\
Andreev, exact & gates, lower & $\Omega(n^{3/2}/\log^3n)$ \cite{KW} & $\Omega(n^{3/2}/\log^2n)$ & Thm.~\ref{thm:lower}\\
Andreev, $\varepsilon$-agree & gates, upper & none stated & $O(\varepsilon^2n^{3/2}\sqrt{\log n})$ & Thm.~\ref{thm:avgupper}\\
Andreev, $\varepsilon$-agree & gates, lower & $\Omega(\varepsilon^3n^{3/2}/\log^3n)$ \cite{KW} & $\Omega(\varepsilon^2n^{3/2}/\log^2n)$ & Thm.~\ref{thm:lower}\\
Andreev, $\varepsilon$-agree & wires, lower & $\Omega(\varepsilon^3n^{5/2}/\log^{7/2}n)$ \cite{KW} & $\Omega(\varepsilon^2n^{5/2}/\log^{5/2}n)$ & Thm.~\ref{thm:lower}\\
some $F\in\mathsf P$, exact & wires, lower & $n^{5/2}/(\log n)^{O(1)}$ \cite{KKL} & -- & \\
Andreev, depth 3 & top fan-in $t$ & $\tilde\Omega(n^{1/2}/s^{1/3})$ \cite{KW} & $\tilde\Omega(n^{3/4}/s^{1/2})$ & Cor.~\ref{cor:fanin}\\
$B_n$, depth 3, exact & gates, lower & $\Omega(n^{3/2}/\log^3n)$ \cite{KW} & $\Omega(n^{3/2}/\log^2n)$ & Thm.~\ref{thm:B}\\
$B_n$, depth 3, exact & wires, lower & $\Omega(n^{5/2}/\log^{7/2}n)$ \cite{KW} & $\Omega(n^{5/2}/\log^{5/2}n)$ & Thm.~\ref{thm:B}\\
$B_n$, $\varepsilon$-agree & gates, lower & none stated & $\Omega(\varepsilon n^{3/2}/\log^2n)$ & Thm.~\ref{thm:B}\\
\midrule
\multicolumn{5}{@{}l}{\emph{Functions in $\mathsf{E}^{\mathsf{NP}}$}}\\
some $F$, exact & gates, lower & $n^{5/2-\delta}$ for every $\delta>0$ \cite{CTW} & -- & \\
some $F$, $\varepsilon$-agree & wires, lower & $n^3/\log^{12}n$ \cite{Nag} (unreviewed) & -- & \\
\bottomrule
\end{tabular}
\end{center}
All bounds are for $\THR\circ\THR$ unless the row says depth~3; $n$ is the
input length; $\varepsilon$-agreement bounds hold for $\varepsilon\ge
C\log n/\sqrt n$; the $B_n$ rows use the parameters of Kane and Williams.
The upper bound $\tfrac34n^2$ shows that no gate lower bound above $n^2$ is
possible for Andreev's function; the $\mathsf E^{\mathsf{NP}}$ bounds therefore
cannot be obtained for it.

\smalltitle{How to read this paper.}
Sections~\ref{sec:upper} and~\ref{sec:avgupper} (upper bounds) and
Section~\ref{sec:lower} (lower bounds) can be read independently.
Each lemma uses only its own local variables plus the few global quantities
in the boxes. All results except Proposition~\ref{prop:modp} are
formalized in Lean (Section~\ref{sec:lean}).

\begin{keybox}
\textbf{Global notation for the upper bounds.}
For $S\subseteq\{1,\dots,k\}$ and $c\in\bits^k$ write $S\cdot c=\sum_{i\in S}c_i$.
\[
X_S=\sum_{c\in\bits^k}(-1)^{S\cdot c+x_c},\qquad
t_S=\sum_{i\in S}s_i .
\]
$X_S$ depends only on the table and lies in $[-N,N]$; it is even.
$t_S$ is the total weight of the blocks in $S$ and lies in $[0,|S|m]$.
\end{keybox}

%-------------------------------------------------------------------
\sect{An exact circuit with \texorpdfstring{$\tfrac34 n^2$}{3/4 n\^{}2} gates}\label{sec:upper}

The circuit computes the integer score $N\cdot(-1)^{x_b}$ and its top gate
tests whether the score is negative. Three lemmas produce the score.

\begin{lemma}[Lookup as a sum of products]\label{lem:lookup}
$\displaystyle\sum_{S\subseteq\{1,\dots,k\}}X_S\,(-1)^{t_S}=N\,(-1)^{x_b}.$
\end{lemma}
\begin{proof}
Since $s_i\equiv b_i\pmod2$, $(-1)^{t_S}=(-1)^{S\cdot b}$. Exchanging sums,
the left side is
\[
\sum_c(-1)^{x_c}\prod_{i=1}^k\bigl(1+(-1)^{c_i+b_i}\bigr),
\]
and the product is $N$ for $c=b$ and $0$ otherwise.
\end{proof}

Let $\varphi(u)=(-1)^u\,u$ for integers $u$.

\begin{lemma}[Removing the product]\label{lem:parity}
If $X$ is even and $t$ is an integer, then $X(-1)^t=\varphi(X+t)-\varphi(t)$.
\end{lemma}
\begin{proof}
$X+t$ has the parity of $t$, so $\varphi(X+t)=(-1)^t(X+t)$.
\end{proof}

\begin{lemma}[Threshold staircase]\label{lem:ladder}
If an integer quantity $L$ always lies in $\{\ell,\dots,h\}$ and $g$ is any
function on the integers, then
\[
g(L)=g(\ell)+\sum_{\theta=\ell+1}^{h}\bigl(g(\theta)-g(\theta-1)\bigr)\ind{L\ge\theta}.
\]
\end{lemma}
\begin{proof}
The brackets are $1$ exactly for $\theta\le L$; the sum telescopes.
\end{proof}

If $L$ is a linear form in the input bits, each $\ind{L\ge\theta}$ is one
threshold gate. Both $X_S+t_S$ and $t_S$ are such forms, because
$(-1)^{x_c}=1-2x_c$:
\begin{equation}\label{eq:form}
X_S+t_S=\sum_c(-1)^{S\cdot c}-2\sum_c(-1)^{S\cdot c}x_c+\sum_{i\in S}\sum_{j=1}^m a_{i,j}.
\end{equation}

\begin{theorem}\label{thm:exact}
Andreev's function has an exact depth-two threshold circuit with
$2N^2+kmN$ bottom gates and one top gate; this is $\tfrac34n^2+1$ gates in
the balanced case. Bottom weights lie in $\{-2,\dots,2\}$; top weights are
odd integers of absolute value below $2n$.
\end{theorem}
\begin{proof}
By Lemmas~\ref{lem:lookup} and~\ref{lem:parity},
\[
N(-1)^{x_b}=\sum_S\bigl(\varphi(X_S+t_S)-\varphi(t_S)\bigr).
\]
Apply Lemma~\ref{lem:ladder} to each term. $X_S+t_S$ ranges over an interval
of length $2N+|S|m$, and $t_S$ over one of length $|S|m$. The score is a
constant plus a weighted sum of these bottom gates, and the top gate outputs
$1$ exactly when the score is negative, i.e.\ when $x_b=1$. Since each index
lies in half of the sets $S$, $\sum_S|S|=kN/2$, and the number of gates is
$\sum_S(2N+2|S|m)=2N^2+kmN$. The top weights are
$\pm(\varphi(\theta)-\varphi(\theta-1))=\pm(-1)^\theta(2\theta-1)$ with
$|\theta|\le n$.
\end{proof}

\smalltitle{Comparison.}
Kane and Williams expand $x_b$ into about $N^2$ parities and build each
parity from its own depth-two subcircuit. Lemma~\ref{lem:lookup} groups the
$N$ parities that share the address character $S$ into the single product
$X_S(-1)^{t_S}$, and Lemma~\ref{lem:parity} removes the product.

\begin{keybox}
\begin{proposition}[Other moduli; not formalized]\label{prop:modp}
Replace parity by $b_i=s_i\bmod p$, with a table indexed by
$\{0,\dots,p-1\}^k$ ($N=p^k$). Then the function has an exact depth-two
threshold circuit with about $pN^2+(p-1)^2kmN$ gates.
\end{proposition}
\emph{Proof sketch.} Expand over the characters of $\mathbb Z_p^k$, and
group the table by the residue $r=S\cdot c\bmod p$ into integer counts
$Z_{S,r}$. Each term is $Z_{S,r}\,g_r(t_S)$ with $g_r(u)=\cos(2\pi(u-r)/p)$
of period $p$, and $pZ\,g(t)=\psi(pZ+t)-\psi(t)$ for $\psi(u)=u\,g(u)$. The
terms $\psi(t_S)$ cancel when summed over $r$. We checked this circuit gate
by gate for $p=3,4,5$ and tables of up to $27$ entries.
\end{keybox}

%-------------------------------------------------------------------
\sect{Agreeing on most inputs with fewer gates}\label{sec:avgupper}

\begin{theorem}\label{thm:avgupper}
For every $\varepsilon\ge C/\sqrt n$ there is a depth-two threshold circuit
with $O(\varepsilon^2n^{3/2}\sqrt{\log n})$ gates that $\varepsilon$-agrees
with Andreev's function. For every $\delta>0$ there is one with
$O(n^{3/2}\sqrt{\log(n/\delta)})$ gates that agrees with it on a $1-\delta$
fraction of inputs.
\end{theorem}

Two ideas save gates.

\smalltitle{Fewer terms.}
Keep only the sets $S\subseteq\{1,\dots,r\}$. The computation in
Lemma~\ref{lem:lookup} now gives
\[
\sum_{S\subseteq\{1,\dots,r\}}X_S(-1)^{t_S}=2^r\sum_{c\in Q}(-1)^{x_c},
\]
where $Q$ is the set of $L=N/2^r$ addresses that agree with $b$ in the first
$r$ coordinates. The sign of this sum is a majority vote of $L$ table
entries, one of which is $x_b$. On a random table the vote equals $x_b$ with
probability at least $\tfrac12+1/(2\sqrt{2L})$: the vote can only be swayed
by $x_b$ when the other $L-1$ entries are balanced, which happens with
probability $\binom{L-1}{L/2}2^{1-L}\ge 1/\sqrt{2L}$.

\smalltitle{Shorter staircases.}
On a random input, Hoeffding's inequality keeps $X_S$ within $\pm\lambda$ of $0$
and $t_S$ within $\pm\lambda$ of $|S|m/2$, where $\lambda=\sqrt{2N\ln(4\cdot2^r/\delta)}$,
for all $2^r$ sets $S$ at once, except with probability $\delta$.
A staircase restricted to such a window of $O(\lambda)$ values computes
$g(L)$ whenever $L$ lies in the window.

\smalltitle{Counting.}
The circuit has $O(2^r\lambda)$ gates. Taking $2^r=N$ gives the second claim.
For the first, take $L\approx1/(32\varepsilon^2)$, so the vote has advantage
about $2\varepsilon$ and $2^r\approx32\varepsilon^2N$, and take $\delta=\varepsilon$.
Then $\log(2^r/\delta)=O(\log n)$. \qed

The lower bounds of Section~\ref{sec:lower} show that both bounds are tight up
to polylogarithmic factors.

%-------------------------------------------------------------------
\sect{Lower bounds}\label{sec:lower}

\begin{theorem}\label{thm:lower}
Let $C$ be a depth-two threshold circuit with $s$ gates and $w$ wires that
$\varepsilon$-agrees with Andreev's function, where
$\varepsilon\ge3(k+2)/\sqrt N$. Then
\[
m\,N\,\varepsilon^2\le 36\,(k+2+\ln(s+1))\cdot s\cdot\sqrt n,
\qquad
m^{3/2}N\varepsilon^2\le 36\,(k+2+\ln(w+1))\cdot w .
\]
In the balanced case this gives $s=\Omega(\varepsilon^2n^{3/2}/\log^2n)$
and $w=\Omega(\varepsilon^2n^{5/2}/\log^{5/2}n)$. Taking $\varepsilon=\tfrac12$
gives the worst-case bounds.
\end{theorem}

\smalltitle{Plan.}
We follow Kane and Williams: fix all but one bit per block at random. The
remaining circuit computes a \emph{random} table lookup on $k$ bits, but
only a few bottom gates still matter. A small circuit cannot correlate with a
random table, so the original circuit cannot either. The improvements over
\cite{KW} are two counting refinements (Lemma~\ref{lem:count} and the
union bound over all survivor counts in Lemma~\ref{lem:avg}).

\begin{keybox}
\textbf{Notation for this section.}
A \emph{restriction} picks one \emph{live} position in every block (the
choice is called $\sigma$) and fixes every other input bit, table included.
The restricted function of the $k$ live bits $u$ is
\[
u\mapsto x_{u\oplus c},
\]
where $c$ is the parity of the fixed bits in each block. A bottom gate
\emph{survives} if it still depends on the live bits; let $U$ be the number
of survivors.
\end{keybox}

\begin{lemma}[Few survivors; Kane--Williams]\label{lem:surv}
For a uniformly random restriction,
$\mathbb E[U]\le 4s\sqrt n/m$. Counting instead the gates that depend on two
or more live bits, the expectation is at most $4w/m^{3/2}$.
\end{lemma}
These are the random restriction lemmas of~\cite[Section~3]{KW}, proved with the Littlewood--Offord
anti-concentration inequality. Their proofs need small repairs (our erratum
lists them); the formalization proves the repaired statements.

\begin{lemma}[Counting restricted circuits]\label{lem:count}
Fix the circuit and $\sigma$. Every restricted circuit with $q$ survivors
belongs to a set of at most
\[
\bigl((2^k+1)M\bigr)^q\,(2^k+1)^{q+k+1}
\]
functions of $u$, where $M=s+1$.
\end{lemma}
\begin{proof}
A surviving gate keeps its original weights on the live bits; only its
threshold depends on the fixed bits. For a fixed weight vector $\beta$, the
gates $\ind{\beta\cdot u\ge\tau}$ are nested as $\tau$ grows, so there are at
most $2^k+1$ of them. Hence each survivor is one of at most $(2^k+1)M$
functions: choose the gate (or a constant), then the threshold. The top gate
is a threshold function of the $q$ survivors and the $k$ live bits; by Chow's
theorem there are at most $(2^k+1)^{q+k+1}$ such functions on a set of $2^k$
points.
\end{proof}
Kane and Williams count all threshold functions of $k$ bits here, which costs
about $k^2$ bits per survivor instead of $O(k+\log s)$. This is where one
logarithm is saved.

\begin{lemma}[Uniform correlation bound]\label{lem:avg}
Let $\alpha=\sqrt{N(k+2+\ln(s+1))}$ and $\beta=(k+2)\sqrt N$. Then the
circuit's agreement with $A$ is at most
\[
\tfrac12+\frac{\alpha\sqrt{\mathbb E[U]}+\beta}{N}+2^{-(k+6)}.
\]
\end{lemma}
\begin{proof}
Fix $\sigma$. Call a fixing \emph{bad} if, for some $q<N$, some function $h$
from Lemma~\ref{lem:count} agrees with $u\mapsto x_{u\oplus c}$ on more than
$N/2+\alpha\sqrt q+\beta$ of the $N$ points $u$.

Given $\sigma$ and the fixed address bits, the restricted table is uniformly
random, so Hoeffding's inequality bounds this for a single $h$ by
$\exp(-2(\alpha\sqrt q+\beta)^2/N)$. Summing over $h$ and $q$ gives total
probability at most $2^{-(k+6)}$ (a geometric series; this is where
$\alpha$ and $\beta$ are chosen).

On a good fixing, the agreement on the subcube is at most
$N/2+\alpha\sqrt U+\beta$. If $U\ge N$ this holds trivially, since
$\alpha\sqrt U\ge N/2$.

Averaging over all restrictions and using
$\mathbb E\sqrt U\le\sqrt{\mathbb E U}$ (Cauchy--Schwarz) proves the lemma.
\end{proof}
Kane and Williams instead apply Markov's inequality to $U$ and then bound a
single survivor count. That loses a factor of $\varepsilon$; bounding all
counts at once and averaging with Cauchy--Schwarz does not.

\begin{proof}[Proof of Theorem~\ref{thm:lower}]
If $\varepsilon\ge3(k+2)/\sqrt N$, the terms $\beta/N$ and $2^{-(k+6)}$ in
Lemma~\ref{lem:avg} add up to at most $2\varepsilon/3$. So
$\varepsilon N/3\le\alpha\sqrt{\mathbb E[U]}$, and squaring gives
$\varepsilon^2N^2\le9\alpha^2\,\mathbb E[U]$.

Insert $\alpha^2=N(k+2+\ln(s+1))$ and Lemma~\ref{lem:surv}. The
wire bound is the same computation with the second part of
Lemma~\ref{lem:surv}, using $w+1$ in place of $s+1$ (each surviving gate is wired to the top gate, so there are at most $w$ of them).
\end{proof}

\smalltitle{Limits of the method.}
After a restriction with one live bit per block, the function
$u\mapsto x_{u\oplus c}$ has a circuit with $N$ gates. Since gates such as
$\ind{t_S\ge\theta}$ survive with probability $\Theta(k/\sqrt N)$, no argument of this form can give more than
$O(n^{3/2}/\log n)$, so Theorem~\ref{thm:lower} is within one logarithm of
the limit. Rectangle methods (discrepancy, sign rank) cannot reach
$n^{2-o(1)}$ either: under every split of the inputs between two players,
$A$ has a protocol with $2k+2$ bits of communication.

%-------------------------------------------------------------------
\sect{Depth three}\label{sec:depth3}

\subsect{Top fan-in for Andreev's function}

A family of depth-two circuits $C_1,\dots,C_t$ is an
\emph{$\varepsilon$-approximate majority} for a function $f$ if on every
input at least a $\tfrac12+\varepsilon$ fraction of the $C_i$ output $f$.

\begin{corollary}\label{cor:fanin}
Every $\varepsilon$-approximate majority for $A$ in which each $C_i$ has at
most $s$ gates satisfies the bounds of Theorem~\ref{thm:lower}. Consequently,
a depth-three circuit whose top gate is a majority of $t$ depth-two
subcircuits with at most $s$ gates each, and which computes $A$, has
\[
t\ \ge\ \tilde\Omega\!\left(\frac{n^{3/4}}{\sqrt s}\right).
\]
\end{corollary}
\begin{proof}
Count, over all inputs and all $i$, the pairs where $C_i$ is correct. Every
input contributes at least $(\tfrac12+\varepsilon)t$ pairs, so some $C_i$ is
correct on a $\tfrac12+\varepsilon$ fraction of inputs. Apply
Theorem~\ref{thm:lower} to that $C_i$.

For a majority of $t$ circuits, more than half are correct on every input,
so $\varepsilon\ge1/(2t)$. As long as $1/(2t)\ge3(k+2)/\sqrt N$,
Theorem~\ref{thm:lower} then gives $s\ge\tilde\Omega(n^{3/2}/t^2)$.
\end{proof}
Kane and Williams obtain $t\ge\tilde\Omega(n^{1/2}/s^{1/3})$. For example, a
depth-three majority circuit with $n^{1.1}$ gates needs top fan-in
$n^{1/5-o(1)}$ rather than $n^{2/15}$.

\subsect{The function \texorpdfstring{$B_n$}{B\_n}}

Andreev's function has depth-three majority circuits of size $O(n)$, so
depth-three lower bounds need a harder function. Kane and Williams use the
following one.

\begin{keybox}
\textbf{Definition.} Let $r_z\in\bits^n$, for $z\in\bits^k$, be a
\emph{$\delta$-biased} list: for every nonzero $y\in\bits^n$,
$\bigl|\sum_z(-1)^{r_z\cdot y}\bigr|\le\delta 2^k$. The input is a seed
$x\in\bits^n$ and $k$ blocks of $\ell$ bits with parities $z$. Then
\[
B(x,a)=r_z\cdot x \bmod 2 .
\]
Kane and Williams take $\ell=n/k$, $k=16\log_2n$ and $\delta=n^{-7}$.
\end{keybox}

The only new ingredient is the following inequality, a version of the
uncorrelation lemma in~\cite[Section~5]{KW} with explicit constants.

\begin{lemma}[Nearly orthogonal vectors]\label{lem:gram}
Let $v_1,\dots,v_t\in\{\pm1\}^M$ satisfy
$|\langle v_i,v_j\rangle|\le\delta M$ for $i\ne j$, and let
$T\in\{\pm1\}^M$. If $|\langle T,v_i\rangle|\ge\gamma M$ for every $i$, then
$t(\gamma^2-\delta)\le1$.
\end{lemma}
\begin{proof}
Let $v=\sum_i\pm v_i$, with signs chosen so that every
$\langle T,\pm v_i\rangle\ge\gamma M$. Then $\langle T,v\rangle\ge t\gamma M$
and $\|v\|^2\le tM+t^2\delta M$. Cauchy--Schwarz gives
$t^2\gamma^2M^2\le M(tM+t^2\delta M)$, which rearranges to the claim.
\end{proof}

For two seeds $x\ne y$, the vectors $u\mapsto(-1)^{B}$ differ by the
character of $x+y$. They are therefore $\delta$-nearly orthogonal, and
Lemma~\ref{lem:gram} shows that a fixed function of $u$ has correlation above
$\gamma$ with at most $1/(\gamma^2-\delta)$ seeds.

\begin{theorem}\label{thm:B}
In the setting of Kane and Williams (formalization in progress, issue \#2812):
\begin{enumerate}
\item every circuit that computes $B$ and consists of a majority gate over
      depth-two threshold circuits has $\Omega(n^{3/2}/\log^2n)$ gates and
      $\Omega(n^{5/2}/\log^{5/2}n)$ wires;
\item every depth-two threshold circuit with $s$ gates agrees with $B$ on at
      most a $\tfrac12+O(s\log^2n/n^{3/2})+n^{-\Omega(1)}$ fraction of
      inputs.
\end{enumerate}
\end{theorem}
\begin{proof}[Proof outline]
We fix the circuit \emph{before} the seed. Kane and Williams fix a seed that
is bad for all small circuits.
\begin{enumerate}
\item For each choice $\sigma$ of live positions, the restricted circuits
      with at most $Q$ survivors, shifted by any $c$, form a set of size
      $2^{O(Q\log n+\log^2 n)}$ (Lemma~\ref{lem:count}). There are only
      $2^{O(\log^2n)}$ choices of $\sigma$. With Lemma~\ref{lem:gram}, the
      fraction of seeds for which some member has correlation above $\gamma$
      is $2^{O(Q\log n+\log^2 n)-n}/(\gamma^2-\delta)$. This is negligible
      when $Q=cn/\log n$.
\item For every fixed seed, $\mathbb E[U]=O(s\log n/\sqrt n)$ by
      Lemma~\ref{lem:surv}.
\item For part 2: the agreement is at most
      $\tfrac12+\gamma+\Pr[\text{seed bad}]+\Pr[U>Q]$, and Markov's
      inequality gives $\Pr[U>Q]\le\mathbb E[U]/Q$. Markov costs nothing
      here, because the correlation of a good restriction is negligible
      rather than about $1/\sqrt n$.
\item For part 1: choose a good seed and a restriction with $U\le Q$. Some
      middle subcircuit must have correlation at least $1/s$ with the
      restricted function, while every one has correlation at most $\gamma$.
      With $\gamma<1/s$ this forces $\mathbb E[U]\gtrsim Q$, i.e.\
      $s\gtrsim n^{3/2}/\log^2 n$. The wire bound is analogous.
\end{enumerate}
\end{proof}
Kane and Williams note that their method cannot show correlation below
$1/\sqrt n$ for circuits of about $n$ gates. Part 2 is consistent with this:
for $s\approx n$ it gives about $\log^2n/\sqrt n$.

%-------------------------------------------------------------------
\sect{Formalization}\label{sec:lean}

All theorems except Proposition~\ref{prop:modp} are proved in Lean~4 in
the Complexitylib library. The proofs check only against the standard
axioms, with no \texttt{sorry}. Files:
\begin{itemize}
\item exact circuit: \texttt{WeightedThreshold/Andreev.lean};
\item average-case circuits: \texttt{AndreevWindowed.lean},
      \texttt{AndreevAverage.lean};
\item restriction lemmas: \texttt{RandomRestriction.lean};
\item counting: \texttt{ThresholdCounting.lean};
\item lower bounds: \texttt{AndreevLowerBound.lean};
\item depth three: see issue~\#2812.
\end{itemize}
The windowed circuits are formalized with the original two-staircase identity
$\varphi(X+t)+\varphi(X-t)=2X(-1)^t$, which changes only constant factors.

\smalltitle{Acknowledgements.}
The results and this note were developed with AI assistance.

\renewcommand{\refname}{\color{ink}References}
\begin{thebibliography}{9}\setlength{\itemsep}{0.1em}
\bibitem{KW} D.~M. Kane and R.~Williams.
\emph{Super-Linear Gate and Super-Quadratic Wire Lower Bounds for Depth-Two
and Depth-Three Threshold Circuits}. STOC 2016.
\href{https://arxiv.org/abs/1511.07860}{arXiv:1511.07860v1}.
\bibitem{KKL} V.~Kabanets, D.~M. Kane, and Z.~Lu.
\emph{A Polynomial Restriction Lemma with Applications}. STOC 2017.
\href{https://eccc.weizmann.ac.il/report/2017/026/}{ECCC TR17-026}.
\bibitem{CTW} L.~Chen, A.~Tal, and Y.~Wang.
\emph{Super-quadratic Lower Bounds for Depth-2 Linear Threshold Circuits}.
STOC 2026. \href{https://eccc.weizmann.ac.il/report/2026/039/}{ECCC TR26-039}.
\bibitem{Nag} D.~Nag. \emph{Almost-Everywhere Near-Cubic Wire Lower Bounds for
$\mathsf{SYM}\circ\THR$ and $\THR\circ\THR$}. Preprint,
\href{https://eccc.weizmann.ac.il/report/2026/167/}{ECCC TR26-167}, 2026.
\end{thebibliography}
\end{document}
